\documentclass{article}
\usepackage[utf8]{inputenc}
\usepackage[T1]{fontenc}
\usepackage{geometry}
\usepackage{graphicx}
\usepackage{float}
\usepackage{hyperref}
\usepackage[english,indonesian]{babel}

\title{Optimasi Pemrosesan Template LaTeX Menggunakan Pendekatan Berbasis AI}
\author{Academic Support Agent}
\date{\today}

\begin{document}
\maketitle

\begin{abstract}

\end{abstract}

\chapter{PENDAHULUAN}


\chapter{TINJAUAN PUSTAKA}


% Original content (commented out):
% \documentclass{article}
\usepackage[utf8]{inputenc}
\usepackage[portuguese]{babel}
\usepackage{draftwatermark}
\usepackage{hyperref}
\usepackage{setspace}
\usepackage{graphicx}
\SetWatermarkAngle{0}
\SetWatermarkText{\includegraphics[width=\paperwidth]{PapelTimbrado.pdf}}

\begin{document}
\begin{center}
    {\bf \huge  1ª Atividade do time \textit{Pyxis}}\\[1.1cm] 
\end{center}
Todas as nossas atividades serão organizadas da seguinte forma: apresentação do objetivo geral e dos objetivos específicos, descrição de uma tarefa em comum e de trabalhos individuais, e o prazo para a entrega. Nossas reuniões semanais devem se basear nas atividades realizadas até a data daquelas. Cada integrante do time terá um tempo para expor os resultados e também haverá a discussão da tarefa em comum. É desejável que esses resultados sejam entregues na forma de texto. A cada 5 atividades, uma apresentação na forma de \textit{slides} deve ser elaborada para uma exposição que deverá ser voltada à equipe Asa Branca.
\begin{doublespace}
\end{doublespace}
Objetivo geral da atividade: reconhecer e descrever, brevemente, a determinação de atitude de um CubeSat.
\begin{doublespace}
\end{doublespace}
Objetivos específicos da atividade:
\begin{itemize}
    \item Assimilar o conceito de GNC (\textit{Guidance, navigation, and control})
    \item Reconhecer os princípios de funcionamento e operação dos sensores estipulados na arquitetura do projeto para a determinação de atitude.
    \item Entender a fusão de sensores.
    \item Compreender o uso do filtro de Kalman para a fusão de sensores e determinação de atitude.
\end{itemize}
\begin{doublespace}
\end{doublespace}
Tarefa em comum da atividade:
\begin{itemize}
    \item Assistir ao vídeo \textit{Introduction to Spacecraft GN\&C - Part 1} - \url{https://goo.gl/bgzY9b}.
    \item Assistir ao vídeo \textit{Introduction to Spacecraft GN\&C - Part 2} - \url{https://goo.gl/RBx6US}.
    \item Pequena leitura \url{https://goo.gl/LV3Ydm}.
\end{itemize}

Lembrando que a atividade em comum também será tema de debate na próxima reunião, assim é recomendada a confecção de uma colinha com os principais aspectos abordados nos vídeos anteriores. 

\begin{doublespace}
\end{doublespace}

Tarefas individuais das atividades:

\begin{doublespace}
\end{doublespace}

Lembrando que são análises curtas, mas que devem carregar a informação necessária para o cumprimento dos objetivos traçados. Nesse primeiro momento, não é preciso descrever profundamente características ou fazer interpretações matemáticas complexas. 

\begin{doublespace}
\end{doublespace}

\begin{itemize}
    \item André Aleixo:
    \begin{itemize}
        \item Elaborar uma lista dos sensores que o projeto utilizará, ressaltando condições de operação (incluindo precisão) e princípios de funcionamento. A arquitetura do projeto deve servir como ponto inicial para seu trabalho.
        \item Artigo recomendado \url{https://goo.gl/R47k1X}
    \end{itemize}
        \item Matheus Roma
    \begin{itemize}
        \item Descrever as possibilidades de fabricação daqueles sensores  que possuem um preço de mercado elevado. Deve ser resaltada a tecnologia aplicada para cada um e a possibilidade de confecção no ambiente da universidade.
        \item Fonte de pesquisa inicial: Trello.
    \end{itemize}
    \item Beatriz Paegle
    \begin{itemize}
        \item Estabelecer uma arquitetura básica para para a determinação de atitude, ressaltando entradas e saídas. Deve ser falado, brevemente, do sistema de coordenadas que o sistema utilizará.
        \item Fontes inicias de pesquisa: \url{https://goo.gl/1NyNxh} e \url{https://goo.gl/SJ66Hi}.
    \end{itemize}
    \item Débora Cardoso
    \begin{itemize}
        \item Descrever a questão da fusão de sensores no contexto da determinação de atitude, ressaltando o erro presente em uma medida e a predição de estados. O algoritmo a ser utilizado é o filtro de Kalman. Você deve arrumar um modo de explicar o algoritmo da forma simples, ressaltando o motivo dele ser utilizado para a fusão de sensores. Haverá uma atividade voltada unicamente para o filtro de Kalman.
        \item Fontes iniciais de pesquisa: \url{https://goo.gl/uY3Hh1}, \url{https://goo.gl/b8Lt4k}, \url{https://goo.gl/6ZaaYR} e \url{https://goo.gl/JV6CYE}.
    \end{itemize}
\end{itemize}

\begin{doublespace}
\end{doublespace}

Os artigos que não estiverem disponíveis gratuitamente para \textit{download} podem ser baixados no \textit{site} \url{http://sci-hub.tw/}. As tarefas individuais possuem algum nível de conexão. Desse modo, seria muito proveitoso que as trocas de informações ocorressem no grupo do \textit{WhatsApp}, assim os laços do time seriam reforçados!

\begin{doublespace}
\end{doublespace}

Prazo de entrga: 03/06/2018.

\end{document}


\end{document}